% Act 1, step 3 (NEW): the sandbox makes sharing safe. One read-only shared state, three children with
% private write layers, inside one policy boundary. Facts: acts/1-introduction/08-runtime-layers.tex.
\begin{frame}[t]{\ev{\tagP\,\tagO}\surf{STEP 3 \textperiodcentered{} SANDBOX}A sandbox lets children share one state safely.}
\vspace{.3cm}
\noindent\begin{tikzpicture}[x=1cm,y=1cm,every node/.style={align=center,font=\fontsize{14}{16.5}\selectfont},
  bl/.style={font=\fontsize{13}{15.5}\selectfont\bfseries,text=Teal,align=center}]
\useasboundingbox (0,.4) rectangle (14.4,-3.75);
% the policy boundary
\draw[Teal,dashed,line width=1.3pt,rounded corners=5pt] (3.85,.05) rectangle (10.55,-3.6);
\node[fill=white,inner xsep=4pt,inner ysep=1pt,bl] at (7.2,.05) {private writes};
% one shared, read-only state
\node[state,minimum width=5.9cm,minimum height=1.0cm] (s) at (7.2,-2.8) {\textbf{Shared state}\enspace{\color{Muted}read-only}};
% three children, each with its own write layer
\foreach \c/\x in {A/5.25,B/7.2,C/9.15}{
  \node[candidate,minimum width=1.5cm,minimum height=1.1cm,font=\fontsize{16}{19}\selectfont\bfseries] (\c) at (\x,-.85) {\c};
  \draw[flow] (\x,-2.3)--(\c.south);}
% undeclared inputs stop at the boundary
\node[bl,text=Rust,anchor=east] at (3.55,-1.3) {undeclared\\inputs hidden};
\draw[Rust,line width=1.2pt,-{Bar[width=4.5mm,line width=1.4pt]}] (1.3,-2.2)--(3.8,-2.2);
% effects leave the boundary
\draw[Gold,line width=1.2pt,-{Latex[length=2.6mm]}] (C.east)--(12.2,-.85);
\node[bl,text=Gold,anchor=west] at (10.85,-1.75) {effects cross\\the boundary};
\end{tikzpicture}
\par\vspace{.1cm}
\claim{A sandbox is a policy; a VM or a container is one way to build it.}
\sources{Diagram: proposed model. Our agent runs used \href{https://docs.docker.com/ai/sandboxes/architecture/}{Docker Sandboxes}: a microVM with a private Docker daemon.}
\qadetail{Proposed model: the shared parent state is immutable (read-only) and named by its digest; each child is that digest plus its own private write layer, so siblings cannot overwrite each other. The boundary has three jobs. Private writes keep alternatives independent until a check accepts one. Hiding undeclared inputs keeps the cache key complete (the step cannot read what the key omits). Effects that cross the boundary, such as a model request, a network call or a push, are not rewound by any restore; the proposal is to record them once (effect journal, idempotency keys), and a restored child must not hold publication authority. Runtime layers are different dimensions rather than interchangeable isolation levels: a conventional VM or microVM runs a guest kernel behind a virtualization boundary (Firecracker uses KVM and a minimal device model); a Linux container shares the host kernel, with namespaces for separate views and cgroups for resource limits; a sandbox is a policy over what execution may access and do, and may be built on a VM, a container, syscall mediation (gVisor) or another mechanism. cgroups alone do not define file access, network policy or syscall permissions; capabilities, seccomp and LSMs are additional mechanisms. Docker Sandboxes, accessed with the sbx CLI, is a microVM-based product with a private Docker daemon and host proxy, not a synonym for an ordinary Docker container. Our factory runs (SELFHOST-2 and SELFHOST-3) ran serially in Docker Sandboxes, not islo, and nothing forked. A sandbox bounds local effects; it does not undo remote ones. No universal startup-time or safety ranking is asserted.}
\note{[0:50] In our design, three children share one saved state. It is read-only and named by its hash. Each child writes only to its private layer, so siblings cannot overwrite each other. Hiding undeclared inputs keeps the cache key complete. Effects, like a network request, cross the boundary. No restore rewinds them, so we propose recording them. This is the Twistlock lesson: a sandbox bounds what a program may do. It is a policy. One way to build it is a VM with its own kernel. Another is a container on a shared kernel. Here are the three steps as code.}
\end{frame}
